% GENERATED FILE. Edit canonical lessons, then run npm run generate:books.
% canonical-source-hash: fnv1a64:568494bf75ec9baf
% canonical-lessons: SA-C469-indriyam, SA-C469-drstih, SA-C469-sparsah, SA-C469-gandhah, SA-C469-jighrati

\chapter{The Senses}
\label{ch:sa-the-senses}
\hlchaptermodality{469}

\begin{chapteropening}
\textbf{By the end of this chapter:} \emph{I can talk about seeing, touching and smelling.}

Complete the last lesson of chapter 469: I can talk about seeing, touching and smelling.
\end{chapteropening}

\section[\texorpdfstring{\sk{इन्द्रियम्} (indriyam)}{इन्द्रियम् (indriyam)}]{\sk{इन्द्रियम्} (indriyam) — a sense; an organ of sense}

\label{lesson:SA-C469-indriyam}



\begin{quote}
Before the new one: say the Sanskrit for reminds, then the Sanskrit for grateful.
\end{quote}

\subsection*{You'll want to know: \sk{इन्द्रियम्}}
\textbf{\sk{इन्द्रियम्}} — \emph{indriyam} — \textquotedblleft{}a sense; an organ of sense\textquotedblright{}.

\begin{quote}
\sk{नरस्य} \sk{पञ्च} \sk{इन्द्रियाणि} \sk{सन्ति।}
\end{quote}

\begin{quote}
\emph{narasya pañca indriyāṇi santi.}
\end{quote}

\begin{quote}
A person has five senses.
\end{quote}

A neuter noun.

\begin{cousinweb}
\textbf{\sk{इन्द्रियम्}} (\emph{indriyam}) is built on \emph{Indra}, the god's name: MW's first sense is \textquotedblleft{}belonging to Indra\textquotedblright{}, then \textquotedblleft{}bodily power\textquotedblright{} and \textquotedblleft{}a sense, an organ of sense\textquotedblright{} (MW).
\end{cousinweb}

\begin{grammarlens}[title={What you've built}]
One more word for describing what you lived and felt.
\end{grammarlens}

\subsection*{Your turn}
\begin{itemize}
\raggedright
  \item \emph{Say it:} \emph{indriyam}
  \item \emph{Say it:} \emph{indriyam}, once more
  \item \emph{Say it:} \emph{narasya pañca indriyāṇi santi.} — A person has five senses.
  \item \emph{Recall:} say the Sanskrit for reminds, then the Sanskrit for grateful, then say \emph{indriyam} again
\end{itemize}

\begin{tcolorbox}[breakable,colback=teal!4,colframe=teal!35!black,title={Before you move on}]
What does \sk{इन्द्रियम्} mean? (\textquotedblleft{}A sense; an organ of sense\textquotedblright{}.) Say it once more.
\end{tcolorbox}
\section[\texorpdfstring{\sk{दृष्टिः} (dṛṣṭiḥ)}{दृष्टिः (dṛṣṭiḥ)}]{\sk{दृष्टिः} (dṛṣṭiḥ) — eyesight}

\label{lesson:SA-C469-drstih}



\begin{quote}
Before the new one: say the Sanskrit for grateful, then the Sanskrit for a sense.
\end{quote}

\subsection*{You'll want to know: \sk{दृष्टिः}}
\textbf{\sk{दृष्टिः}} — \emph{dṛṣṭiḥ} — \textquotedblleft{}eyesight\textquotedblright{}.

\begin{quote}
\sk{वृद्धस्य} \sk{दृष्टिः} \sk{दुर्बला} \sk{अस्ति।}
\end{quote}

\begin{quote}
\emph{vṛddhasya dṛṣṭiḥ durbalā asti.}
\end{quote}

\begin{quote}
The old man's eyesight is weak.
\end{quote}

A feminine noun.

\begin{cousinweb}
\textbf{\sk{दृष्टिः}} (\emph{dṛṣṭiḥ}) comes from the root √dṛś, \textquotedblleft{}to see\textquotedblright{}, the root of \textbf{\sk{दर्शयति}} and \textbf{\sk{दृष्टवान्}}; MW gives \textquotedblleft{}seeing\textquotedblright{} and \textquotedblleft{}sight, the faculty of seeing\textquotedblright{} (MW).
\end{cousinweb}

\begin{grammarlens}[title={What you've built}]
One more word for describing what you lived and felt.
\end{grammarlens}

\subsection*{Your turn}
\begin{itemize}
\raggedright
  \item \emph{Say it:} \emph{dṛṣṭiḥ}
  \item \emph{Say it:} \emph{dṛṣṭiḥ}, once more
  \item \emph{Say it:} \emph{vṛddhasya dṛṣṭiḥ durbalā asti.} — The old man's eyesight is weak.
  \item \emph{Recall:} say the Sanskrit for grateful, then the Sanskrit for a sense, then say \emph{dṛṣṭiḥ} again
\end{itemize}

\begin{tcolorbox}[breakable,colback=teal!4,colframe=teal!35!black,title={Before you move on}]
What does \sk{दृष्टिः} mean? (\textquotedblleft{}Eyesight\textquotedblright{}.) Say it once more.
\end{tcolorbox}
\section[\texorpdfstring{\sk{स्पर्शः} (sparśaḥ)}{स्पर्शः (sparśaḥ)}]{\sk{स्पर्शः} (sparśaḥ) — touch}

\label{lesson:SA-C469-sparsah}



\begin{quote}
Before the new one: say the Sanskrit for a sense, then the Sanskrit for eyesight.
\end{quote}

\subsection*{You'll want to know: \sk{स्पर्शः}}
\textbf{\sk{स्पर्शः}} — \emph{sparśaḥ} — \textquotedblleft{}touch\textquotedblright{}.

\begin{quote}
\sk{मातुः} \sk{स्पर्शः} \sk{कोमलः} \sk{भवति।}
\end{quote}

\begin{quote}
\emph{mātuḥ sparśaḥ komalaḥ bhavati.}
\end{quote}

\begin{quote}
A mother's touch is gentle.
\end{quote}

A masculine noun.

\begin{cousinweb}
\textbf{\sk{स्पर्शः}} (\emph{sparśaḥ}) comes from the root √spṛś, \textquotedblleft{}to touch\textquotedblright{}, the root of \textbf{\sk{स्पृशति}}; MW gives \textquotedblleft{}touch, the sense of touch\textquotedblright{} (MW).
\end{cousinweb}

\begin{grammarlens}[title={What you've built}]
One more word for describing what you lived and felt.
\end{grammarlens}

\subsection*{Your turn}
\begin{itemize}
\raggedright
  \item \emph{Say it:} \emph{sparśaḥ}
  \item \emph{Say it:} \emph{sparśaḥ}, once more
  \item \emph{Say it:} \emph{mātuḥ sparśaḥ komalaḥ bhavati.} — A mother's touch is gentle.
  \item \emph{Recall:} say the Sanskrit for a sense, then the Sanskrit for eyesight, then say \emph{sparśaḥ} again
\end{itemize}

\begin{tcolorbox}[breakable,colback=teal!4,colframe=teal!35!black,title={Before you move on}]
What does \sk{स्पर्शः} mean? (\textquotedblleft{}Touch\textquotedblright{}.) Say it once more.
\end{tcolorbox}
\section[\texorpdfstring{\sk{गन्धः} (gandhaḥ)}{गन्धः (gandhaḥ)}]{\sk{गन्धः} (gandhaḥ) — a smell}

\label{lesson:SA-C469-gandhah}



\begin{quote}
Before the new one: say the Sanskrit for eyesight, then the Sanskrit for touch.
\end{quote}

\subsection*{You'll want to know: \sk{गन्धः}}
\textbf{\sk{गन्धः}} — \emph{gandhaḥ} — \textquotedblleft{}a smell\textquotedblright{}.

\begin{quote}
\sk{पुष्पस्य} \sk{गन्धः} \sk{मधुरः} \sk{अस्ति।}
\end{quote}

\begin{quote}
\emph{puṣpasya gandhaḥ madhuraḥ asti.}
\end{quote}

\begin{quote}
The flower's smell is sweet.
\end{quote}

A masculine noun.

\begin{cousinweb}
\textbf{\sk{गन्धः}} (\emph{gandhaḥ}) MW gives \textquotedblleft{}smell, odour\textquotedblright{} and lists nine kinds of smell; MW gives no derivation, so its origin is uncertain (MW).
\end{cousinweb}

\begin{grammarlens}[title={What you've built}]
One more word for describing what you lived and felt.
\end{grammarlens}

\subsection*{Your turn}
\begin{itemize}
\raggedright
  \item \emph{Say it:} \emph{gandhaḥ}
  \item \emph{Say it:} \emph{gandhaḥ}, once more
  \item \emph{Say it:} \emph{puṣpasya gandhaḥ madhuraḥ asti.} — The flower's smell is sweet.
  \item \emph{Recall:} say the Sanskrit for eyesight, then the Sanskrit for touch, then say \emph{gandhaḥ} again
\end{itemize}

\begin{tcolorbox}[breakable,colback=teal!4,colframe=teal!35!black,title={Before you move on}]
What does \sk{गन्धः} mean? (\textquotedblleft{}A smell\textquotedblright{}.) Say it once more.
\end{tcolorbox}
\section[\texorpdfstring{\sk{जिघ्रति} (jighrati)}{जिघ्रति (jighrati)}]{\sk{जिघ्रति} (jighrati) — smells}

\label{lesson:SA-C469-jighrati}



\begin{quote}
Before the new one: say the Sanskrit for touch, then the Sanskrit for a smell.
\end{quote}

\subsection*{You'll want to know: \sk{जिघ्रति}}
\textbf{\sk{जिघ्रति}} — \emph{jighrati} — \textquotedblleft{}smells\textquotedblright{}.

\begin{quote}
\sk{बालिका} \sk{पुष्पं} \sk{जिघ्रति।}
\end{quote}

\begin{quote}
\emph{bālikā puṣpaṃ jighrati.}
\end{quote}

\begin{quote}
The girl smells the flower.
\end{quote}

Past (she …): \sk{सा} \sk{पुष्पं} \sk{घ्रातवती}, \emph{sā puṣpaṃ ghrātavatī}.

\begin{cousinweb}
\textbf{\sk{जिघ्रति}} (\emph{jighrati}) comes from the root √ghrā, \textquotedblleft{}to smell\textquotedblright{}, whose present stem doubles the root (\emph{ji-ghr-a}) (MW).
\end{cousinweb}

\begin{grammarlens}[title={What you've built}]
One more word for describing what you lived and felt.
\end{grammarlens}

\subsection*{Your turn}
\begin{itemize}
\raggedright
  \item \emph{Say it:} \emph{jighrati}
  \item \emph{Say it:} \emph{jighrati}, once more
  \item \emph{Say it:} \emph{bālikā puṣpaṃ jighrati.} — The girl smells the flower.
  \item \emph{Recall:} say the Sanskrit for touch, then the Sanskrit for a smell, then say \emph{jighrati} again
\end{itemize}

\begin{tcolorbox}[breakable,colback=teal!4,colframe=teal!35!black,title={Before you move on}]
What does \sk{जिघ्रति} mean? (\textquotedblleft{}Smells\textquotedblright{}.) Say it once more.
\end{tcolorbox}
